\documentclass[aspectratio=169,12pt]{beamer}
\usepackage{fontspec}
\usepackage[portuguese]{babel}
\usepackage{xcolor,listings,tikz,booktabs,amsmath,amssymb,graphicx,multicol,array,colortbl}
\usetikzlibrary{arrows.meta,positioning,shapes.geometric,calc}

\definecolor{navy}{HTML}{0D1B3E}
\definecolor{gold}{HTML}{F5A623}
\definecolor{qgreen}{HTML}{1E8B4C}
\definecolor{qblue}{HTML}{1A6EAE}
\definecolor{codebg}{HTML}{EEF3F8}

\usetheme{default}
\useinnertheme{circles}
\setbeamertemplate{navigation symbols}{}
\setbeamertemplate{footline}{%
  \hfill{\color{gray!60}\scriptsize\insertframenumber/\inserttotalframenumber}\hspace{.6em}\vspace{.4em}}

\setbeamercolor{frametitle}{bg=navy,fg=white}
\setbeamercolor{structure}{fg=navy}
\setbeamercolor{alerted text}{fg=gold}
\setbeamercolor{item}{fg=gold}
\setbeamercolor{subitem}{fg=navy}
\setbeamercolor{block title}{bg=qblue,fg=white}
\setbeamercolor{block body}{bg=qblue!8,fg=black}
\setbeamercolor{block title alerted}{bg=gold,fg=navy}
\setbeamercolor{block body alerted}{bg=gold!15,fg=navy}
\setbeamercolor{block title example}{bg=qgreen,fg=white}
\setbeamercolor{block body example}{bg=qgreen!10,fg=black}
\setbeamerfont{frametitle}{size=\large,series=\bfseries}
\setbeamerfont{block title}{series=\bfseries}

\setbeamertemplate{title page}{%
  \begin{tikzpicture}[remember picture,overlay]
    \fill[navy](current page.north west)rectangle(current page.south east);
    \fill[gold]([yshift=.7cm]current page.south west)rectangle(current page.south east);
    \node[anchor=south,yshift=.73cm,navy,font=\tiny\bfseries]at(current page.south)
      {INTENSIVÃO QUANT AI \textbullet{} IMPACT UFSCAR};
  \end{tikzpicture}
  \vspace{.65cm}
  \begin{center}
    {\color{gold}\scriptsize\bfseries INTENSIVÃO QUANT AI \textbullet{} IMPACT UFSCAR}\\[.4cm]
    {\color{white}\LARGE\bfseries\inserttitle}\\[.25cm]
    {\color{gold!85}\normalsize\insertsubtitle}\\[.5cm]
    {\color{white!70}\small\insertauthor}\\[.1cm]
    {\color{white!50}\footnotesize\insertdate}
  \end{center}}

\lstdefinestyle{python}{language=Python,
  basicstyle=\ttfamily\scriptsize\color{qblue},
  keywordstyle=\color{navy}\bfseries,
  stringstyle=\color{qgreen!80!black},
  commentstyle=\color{gray!70}\itshape,
  backgroundcolor=\color{codebg},
  frame=l,framerule=2pt,rulecolor=\color{gold},
  xleftmargin=1em,breaklines=true,showstringspaces=false,tabsize=4,
  extendedchars=false,inputencoding=ascii}
\lstset{style=python}

\author{Felipe Maldonado\\{\scriptsize VP Diretoria Quant~\textbullet{}~Impact UFSCAR}}
\date{Intensivão Quant AI 2026}
\title{Aula 06 --- Alocação}
\subtitle{Alocação Clássica vs Otimizada (Markowitz)}

\begin{document}

\begin{frame}[plain]
\titlepage
\end{frame}

\begin{frame}[fragile]{Nesta Aula}
\begin{columns}[T]
  \column{0.5\textwidth}
\begin{block}{Teoria (20 min)}
\begin{itemize}
  \item Teoria Moderna de Portfólio (Markowitz 1952)
  \item O trade-off risco-retorno e fronteira eficiente
  \item Por que a otimização de Markowitz clássica é instável na prática
  \item A importância de restrições rígidas (long-only, caps)
\end{itemize}
\end{block}
  \column{0.47\textwidth}
\begin{block}{Código (40 min)}
\begin{itemize}
  \item \texttt{pesos\_equal\_weight()} — robustez do benchmark 1/N
  \item \texttt{pesos\_markowitz()} — otimização com scipy.optimize
  \item Configurar função objetivo (max Sharpe) e restrições
  \item Comparação de turnover e Sharpe de Markowitz vs. EW
  \item Salvar \texttt{pesos\_markowitz.parquet}
\end{itemize}
\end{block}
\end{columns}
\end{frame}

\begin{frame}[fragile]{Equal-Weight vs Markowitz}
\begin{columns}[T]
  \column{0.5\textwidth}
\begin{block}{DeMiguel et al.\ (2009)}

\textit{Optimal versus naive diversification: how inefficient is the 1/N portfolio strategy?}
\vspace{.3em}
\begin{itemize}
  \item O portfólio 1/N (Equal-Weight) é um benchmark fortíssimo
  \item Otimização clássica sofre com erro de estimação dos parâmetros ($\boldsymbol{\mu}, \boldsymbol{\Sigma}$)
  \item Sem restrições, Markowitz gera pesos extremos e instáveis out-of-sample
  \item Solução prática: impor restrições long-only e limites máximos por ativo (caps)
\end{itemize}
\end{block}\vspace{.3em}
\begin{alertblock}{Restrições no scipy.optimize}
Impedir short selling ($w_i \in [0, 0.20]$) e forçar a soma dos pesos a ser exatamente $1$. Isso estabiliza os pesos.
\end{alertblock}
  \column{0.47\textwidth}

\textbf{Problema de Markowitz (Sharpe Máx.):}
\[
\max_{\mathbf{w}}\; \frac{\mathbf{w}^\top\boldsymbol{\mu} - r_f}{\sqrt{\mathbf{w}^\top\boldsymbol{\Sigma}\mathbf{w}}}
\]
\vspace{.1em}
\textbf{Com as Restrições:}
\[
\sum_{i=1}^N w_i = 1 \quad \text{e} \quad 0 \leq w_i \leq 0{,}20
\]

\vspace{.3em}
\begin{exampleblock}{Estimativa de Risco}
Usamos a matriz de covariância histórica dos retornos mensais do portfólio selecionado para a otimização.
\end{exampleblock}
\end{columns}
\end{frame}

\begin{frame}[fragile]{O que Vamos Construir}
\begin{columns}[T]
  \column{0.52\textwidth}
\begin{block}{Funções a implementar}
\begin{itemize}
  \item \texttt{pesos\_equal\_weight(sinal)} $\to$ DataFrame de pesos iguais
  \item \texttt{pesos\_markowitz(sinal, ret\_mensais)} $\to$ pesos otimizados (Sharpe Máx.)
  \item \texttt{calcular\_turnover(pesos)} $\to$ Series de turnover mensal
\end{itemize}
\end{block}
  \column{0.45\textwidth}
\begin{block}{Entradas (parquets)}
\begin{itemize}
  \item \texttt{sinal\_v1.parquet}
  \item \texttt{retornos\_mensais\_limpo.parquet}
\end{itemize}
\end{block}
\vspace{.4em}
\begin{exampleblock}{Saídas (parquets)}
\begin{itemize}
  \item \texttt{pesos\_markowitz.parquet}
  \item Tabela comparativa de Sharpe e turnover: Markowitz vs EW
\end{itemize}
\end{exampleblock}
\end{columns}
\end{frame}

\begin{frame}[fragile]{Referências}
\begin{multicols}{2}
\tiny
\begin{itemize}
  \item \textbf{Markowitz, H.} (1952). Portfolio Selection. \textit{Journal of Finance}, 7(1), 77--91.
  \item \textbf{DeMiguel, V. et al.} (2009). Optimal versus naive diversification. \textit{Review of Financial Studies}, 22(5), 1915--1953.
  \item \textbf{Ledoit, O. \& Wolf, M.} (2004). Honey, I shrunk the covariance matrix. \textit{JPM}.
  \item \textbf{Michaud, R.} (1989). The Markowitz optimization enigma: is optimized optimal? \textit{FAJ}, 45(1), 31--42.
\end{itemize}
\end{multicols}
\end{frame}


\end{document}
